% !TEX root = ../../Main.tex
\fancychapter{Conclusion}
\label{Chapter:Conclusion}

The aim of this work was to find out whether one deep reinforcement learning design gives a profitable policy with controlled risk on all seven major currency pairs. The design was trained separately for each instrument and was otherwise not changed. The evaluation charged the same costs that a real account charges.

The answer is yes, but with conditions. These conditions are as much a part of the result as the answer itself.

% !TEX root = ../../Main.tex
\section{Summary}
\label{Section:Summary}

This work built three things and answered one question.

The first is a trading framework. In it, one strategy implementation runs in three places: an offline backtest over stored quotes, the cTrader backtesting tab, and a live account. Its execution engine reproduces the platform's own results byte for byte on six reference configurations. It also goes further than the platform. It steps through every tick instead of only the bar closes, and it charges the spread recorded at the moment of each trade. The framework reads from a market database of $1.66$ billion quotes for the seven pairs, which takes up $295$ GB. Every bar series used in this work is built from this database.

The second is an agent. It observes thirty-eight dimensionless values, built from sixteen technical indicators. It outputs one continuous action that sets both the direction and the size of its position at the same time. Six changes to the published algorithm were needed to make it train on this problem, and two of them matter most. The first is a scale-sensitive reward. A scale-invariant reward gives no gradient towards holding a position large enough to matter. The second is a penalty on the actor's pre-activation. Without it, the policy saturates at $+1$ or $-1$ and collapses to a constant action.

The third is a selection procedure. First, gates reject candidates with degenerate policies, such as those that hardly trade or keep to one side of the market. Next, each candidate is tested against a null built from itself, by rotating the model's own sequence of positions in time. Last, the candidates are ranked by a composite score in which return is one of nine parts. The procedure was not designed to be elegant. It was built because ranking on return does not work, and \Cref{Section:StatisticalScreening} shows this with measured results.

The answer to the question is as follows. All seven agents are profitable over the eleven-year window, from each of the five starting balances. Their total returns range from $+15.7\%$ to $+226.6\%$ net of the spread and commission of a swap-free raw-spread account. Six of the seven have a positive alpha when their return is decomposed against their own market. Three of these six do this with almost no market exposure. The seventh, USD/JPY, only follows its market. It is reported as a failure, together with a diagnosis of its cause. An equal-weight portfolio of the seven lowers the maximum drawdown from an average of $34.5\%$ to $15.9\%$ and raises the Sharpe ratio to $0.632$. This happens because the agents are close to uncorrelated with one another.

Two findings lie behind these numbers. They are more likely than the numbers themselves to carry over to other problems. First, the agents really use the continuous action. The mean absolute action ranges from $0.09$ to $0.98$ across the pairs, and saturation is below one percent on every pair. So the agents learned to size their positions, and they do not just switch between the two extremes. Second, win rates lie between $49.8\%$ and $55.0\%$, and the strongest model in the set has the lowest win rate. So the profit does not come from being right more often. It comes from what the agents do with a position once it is open. A continuous action makes this possible and a discrete action does not.

% !TEX root = ../../Main.tex
\section{Limitations}
\label{Section:Limitations}

The most important limitation comes first, because it puts a limit on everything else. The performance figures in \Cref{Chapter:ExperimentalResults} are measured over a window that includes the years the agents were trained on. The procedure does keep one year out of training and out of the choice among an agent's own episodes. But the choice of which trained agent to publish was made on the full window, so the headline returns are an in-sample result. On the one year that no part of the procedure saw, the models lost money. The portfolio returned $-9.8\%$ in 2025, and five of the seven pairs were negative. That year had a broad dollar reversal, and the training decade had none. The one model that held up is the least directional model in the set. However, a single year is a small sample, and no stronger conclusion should be drawn from it in either direction.

There is a second way in which the selection is not corrected. Each published model was chosen from between $128$ and $256$ candidates. With that many draws, a permutation $p$-value near $0.05$ is close to what chance alone produces. For this reason, the screening is reported as a filter that was shown to reject bad models. It is not reported as a claim that the models that passed are statistically significant. Under a strict correction, only USD/CAD and GBP/USD would clearly pass.

Three smaller limitations should also be recorded. First, the position size was chosen for each pair after the results were seen. It scales every trade and does not change the policy that makes the trades. Still, it is a parameter fitted on the same data as everything else. Second, training cannot be reproduced across environments. It is bit-exact within one installation on a single thread, but a different version of a numerical library changes the outcome. So the delivered models are archived artifacts whose replay has been verified. They are not a recipe that can produce the same models again. Third, two weeks of January 2016 are missing from all seven pairs, for reasons on the broker's side. The gap is the same for all pairs and lies outside the evaluation window, so it does not bias any comparison in this work.

Finally, the tests cover a narrower scope than the method itself. They use seven instruments from one asset class, one eleven-year period, one decision frequency and one cost model. This is only one case out of many possible ones. Nothing in this work shows that the design carries over to other markets, other periods or other holding horizons. The USD/JPY result shows that it does not carry over to every instrument, even among the seven.

% !TEX root = ../../Main.tex
\section{Future Work}
\label{Section:FutureWork}

The limitations above set the order of the future work. The most useful next step is the one that would most change the value of the results.

The first step is a true held-out evaluation. Training would end before the reporting period starts. The choice among candidates would use only data from before that point. The reported figures would then be an out-of-sample estimate. They would no longer describe the same window on which the models were chosen. Of all possible changes, this one would do the most to make the claims stronger. The framework already supports it. It needs a campaign run under this stricter protocol, and it needs no new features in the framework.

The second step follows from the first. The procedure used here splits the window only once. A walk-forward scheme with several folds would train each fold on the data before it and evaluate it on the data after it. This would give an out-of-sample record for the whole period, and not just for one year. The learned weights could also be carried from one fold into the next, instead of starting each fold from a new initialisation. This would test whether an agent adapts as the market changes. The framework supports this, but it was not used in this campaign.

The third step is to report how much the results vary between training runs. The path-robustness protocol used here changes the starting balance. This measures how sensitive a result is to the entry path, but not how sensitive it is to the training itself. The next step would report the distribution of outcomes across the whole pool of candidates, and where the published model falls within it. A reader could then see how much of a result comes from the method and how much comes from chance.

Apart from the evaluation, what the agents did suggests three further directions. The portfolio in \Cref{Section:Portfolio} gives the seven agents equal weights. Weighting them by risk, or by their correlation structure, should do better. The seven agents are almost independent of each other, which suggests that the gain would be real. The observation holds thirty-eight values from a single point in time. The natural next step is a recurrent or attention-based encoder that reads a window of past data directly. This was part of the original plan for this work. Last, the framework already runs the same strategy on a live account. A live deployment over a long period would test whether the execution assumptions hold when the orders are real. No backtest can test this.

One improvement to the framework itself should also be noted. The execution engine charges the spread correctly, by filling orders at the ask and at the bid. However, the exported trade record does not show this cost as a separate field. To recover it, the indirect calculation of \Cref{Equation:SpreadRecovery} is needed. If every trade recorded the spread cost next to its commission and swap, the cost decomposition in \Cref{Tables:Costs} could be read directly and would not have to be inferred. It could also be reported for each position, and not only in total.

The USD/JPY result needs a separate note. The current explanation of the failure on this pair is that the price window is monotonic, which means the price moves in one direction only. Such a window gives no gradient towards changing the size of the exposure. It has not been shown that this is the whole explanation. The pair should be studied again once the evaluation protocol above is in place.

Some reported returns suggest that deep reinforcement learning for currency trading is a solved problem. Its out-of-sample record suggests that it leads nowhere. Neither view is correct. This work shows that the main difficulty is not in the learning algorithm. It is in the environment the agent is trained in and in the costs it is charged. Most of all, it is in a careful and strict way of deciding which of many trained models to trust. Six of the seven agents produced returns that still hold after a decomposition into policy and market exposure. For all seven pairs, the candidate that looked best on return did not pass this test.